% Шаблон документа для курса "Системы научных публикаций"
% Семинар 14. LaTeX. Вставка и настройка библиографической информации
\documentclass[a4paper,12pt]{article}

%%% Работа с русским языком
\usepackage{cmap}
\usepackage{mathtext}
\usepackage[T2A]{fontenc}
\usepackage[utf8]{inputenc}
\usepackage[english,russian]{babel}
\usepackage{indentfirst}

\usepackage{extsizes}
\usepackage{geometry}
\geometry{top=25mm}
\geometry{bottom=35mm}
\geometry{left=35mm}
\geometry{right=20mm}

\usepackage{setspace}
\usepackage{lastpage}
\usepackage{soul}
\usepackage{multicol}

\usepackage{hyperref}
\usepackage[usenames,dvipsnames,svgnames,table,rgb]{xcolor}
\hypersetup{
  unicode=true,
  pdftitle={Заголовок},
  pdfauthor={Автор},
  pdfsubject={Тема},
  pdfcreator={Создатель},
  pdfproducer={Производитель},
  pdfkeywords={keyword1} {key2} {key3},
  colorlinks=true,
  linkcolor=black,
  citecolor=black,
  filecolor=magenta,
  urlcolor=blue
}

%%% Дополнительная работа с математикой
\usepackage{amsfonts,amssymb,amsthm,mathtools}
\usepackage{amsmath}
\usepackage{icomma}

\usepackage{euscript}
\usepackage{mathrsfs}

\DeclareMathOperator{\sgn}{\mathop{sgn}}

\newcommand*{\hm}[1]{#1\nobreak\discretionary{}
{\hbox{$\mathsurround=0pt #1$}}{}}

%%% Работа с таблицами
\usepackage{array,tabularx,tabulary,booktabs}
\usepackage{longtable}
\usepackage{multirow}
\usepackage{wrapfig}

%%% Работа с рисунками
\usepackage{graphicx}
\graphicspath{{images/}{images2/}}
\setlength\fboxsep{3pt}
\setlength\fboxrule{1pt}

\usepackage{cite}
\usepackage{csquotes}

%%% Заголовок
\author{Соин Андрей Дмитриевич}
\title{\LaTeX. Вставка и настройка библиографической информации}
\date{\today}

\begin{document}

\maketitle

\tableofcontents

\section{Введение}
Пример ссылок на источники:
- на один источник: \cite{lvovsky};
- на два источника не по порядку: \cite{kim,filimonov2012};
- на три источника подряд: \cite{kim,philips,filimonov2012}.

Ссылка на источник с меткой: \cite{lvovsky}.

\section*{thebibliography}

\addcontentsline{toc}{section}{Список использованных источников}

\renewcommand{\refname}{Список использованных источников}

\makeatletter
\renewcommand{\@biblabel}[1]{#1.}
\makeatother

\begin{thebibliography}{9}

  \bibitem{lvovsky} Львовский С. Набор и верстка в системе \LaTeX. — Litres, 2022.

  \bibitem{kim} Ким Д. П. Теория автоматического управления. — Физматлит, 2003.

  \bibitem{philips} Филипс Ч. Системы управления с обратной связью. — Лаборатория базовых знаний, 2001.

  \bibitem{filimonov2012} Филимонов А. Б., Филимонов Н. Б. О проблематике синтеза координирующих систем автоматического управления // Известия Южного федерального университета. Технические науки. — 2012. — Т. 128. — № 3. — С. 172–180.

  \bibitem{zhang2023} Zhang W. et al. Converted state equation Kalman filter for nonlinear maneuvering target tracking // Signal Processing. — 2023. — Vol. 202. — P. 108–741.

\end{thebibliography}

\end{document}